\documentclass[a4paper,11pt]{article}
\usepackage[T1]{fontenc}
\usepackage[utf8]{inputenc}
\usepackage[italian]{babel}
\usepackage{amsmath,amssymb}
\usepackage{graphicx}
\usepackage{geometry}
\usepackage{xcolor}
\usepackage{listings}
\geometry{margin=2.5cm}

\lstdefinelanguage{MiniZinc}{
  morekeywords={int,set,of,array,var,constraint,include,solve,satisfy,output,show},
  sensitive=true,
  morecomment=[l]{\%},
  morestring=[b]"
}

\lstset{
  language=MiniZinc,
  basicstyle=\ttfamily\small,
  keywordstyle=\color{blue}\bfseries,
  commentstyle=\color{teal!70!black}\itshape,
  stringstyle=\color{orange!70!black},
  numbers=left,
  numberstyle=\tiny\color{gray},
  stepnumber=1,
  numbersep=8pt,
  showstringspaces=false,
  frame=single,
  breaklines=true
}

\title{FOL: Studenti Ansiosi}
\date{}

\begin{document}
\maketitle

\section*{Problema: Studenti Ansiosi}
Modellare mediante una formula $\phi$ in logica del primo ordine la seguente affermazione:

\textit{Nessuno studente supererà un qualunque esame se è ansioso o non ne studia il programma}

\subsection*{1.1 Modellazione}


\paragraph{Predicati}
\begin{itemize}
  \item $\mathrm{Studente}(x)$: $x$ è uno studente
  \item $\mathrm{Esame}(y)$: $y$ è un esame
  \item $\mathrm{Ansioso}(x)$: $x$ è ansioso
  \item $\mathrm{StudiaProgramma}(x, y)$: $x$ studia il programma dell'esame $y$
  \item $\mathrm{Supera}(x, y)$: $x$ supera l'esame $y$
\end{itemize}

\paragraph{Formula}

La formula che modella l'affermazione è:
\[
\phi = \forall x \, \forall y \,
\left(
  \left(
    \mathrm{Studente}(x) \land \mathrm{Esame}(y) \land
    \left( \mathrm{Ansioso}(x) \lor \neg \mathrm{StudiaProgramma}(x, y) \right)
  \right)
\Rightarrow
\neg \mathrm{Supera}(x, y)
\right)
\]

\newpage

\subsection*{1.2 Modello $M$ di $\phi$}

Consideriamo l'interpretazione $M$ con:
\begin{itemize}
  \item Dominio = $\{\mathrm{pino}, \mathrm{AI}\}$
\end{itemize}

\paragraph{Interpretazione}
\begin{itemize}
  \item $\mathrm{Studente}^M = \varnothing$
  \item $\mathrm{Esame}^M = \{\mathrm{AI}\}$
  \item $\mathrm{Ansioso}^M = \{\mathrm{pino}\}$
  \item $\mathrm{StudiaProgramma}^M = \varnothing$
  \item $\mathrm{Supera}^M = \varnothing$
\end{itemize}

\paragraph{Verifica che $M \models \phi$}

Poiché non esiste alcun elemento che soddisfi il predicato $\mathrm{Studente}$, l'antecedente della formula
\[
\mathrm{Studente}(x) \land \mathrm{Esame}(y) \land \left(\mathrm{Ansioso}(x) \lor \neg \mathrm{StudiaProgramma}(x,y)\right)
\]
è sempre falso per ogni scelta di $x$ e $y$. Di conseguenza, l'implicazione presente in $\phi$ risulta vera per ogni scelta di $x$ e $y$.

Pertanto, $M \models \phi$.

\newpage

\subsection*{1.3 Non-modello $I$ di $\phi$}

Consideriamo l'interpretazione $I$ con:
\begin{itemize}
  \item Dominio = $\{\mathrm{pino}, \mathrm{AI}\}$
\end{itemize}

\paragraph{Interpretazione}
\begin{itemize}
  \item $\mathrm{Studente}^I = \{\mathrm{pino}\}$
  \item $\mathrm{Esame}^I = \{\mathrm{AI}\}$
  \item $\mathrm{Ansioso}^I = \{\mathrm{pino}\}$
  \item $\mathrm{StudiaProgramma}^I = \varnothing$
  \item $\mathrm{Supera}^I = \{(\mathrm{pino}, \mathrm{AI})\}$
\end{itemize}

\paragraph{Verifica che $I \not\models \phi$}

Per $x=\mathrm{pino}$ e $y=\mathrm{AI}$ si ha:
\[
\left(
  \mathrm{Studente}(\mathrm{pino})
  \land
  \mathrm{Esame}(\mathrm{AI})
  \land
  \left(
    \mathrm{Ansioso}(\mathrm{pino})
    \lor
    \neg \mathrm{StudiaProgramma}(\mathrm{pino},\mathrm{AI})
  \right)
\right)
\Rightarrow
\neg \mathrm{Supera}(\mathrm{pino},\mathrm{AI})
\]
L'antecedente dell'implicazione è vero, ma il conseguente
\[
\neg \mathrm{Supera}(\mathrm{pino},\mathrm{AI})
\]
è falso poiché $(\mathrm{pino}, \mathrm{AI}) \in \mathrm{Supera}^I$. Pertanto l'implicazione è falsa.

Pertanto, $I \not\models \phi$.


\end{document}
